\documentclass[11pt,a4paper]{article}

% Appunti personali di lettura: non sono un testo da presentare.
\usepackage[utf8]{inputenc}
\usepackage[T1]{fontenc}
\usepackage[italian]{babel}
\usepackage{amsmath,amssymb,mathtools}
\usepackage[a4paper,margin=2.45cm]{geometry}
\usepackage{enumitem}
\usepackage{booktabs}
\usepackage{xcolor}
\usepackage{tikz}
\usetikzlibrary{arrows.meta,positioning,calc,patterns}
\usepackage{hyperref}

\title{Appunti di lettura\\\large Hua e Joe (2011), ``Tail order and intermediate tail dependence''}
\author{Appunti personali di Massimo}
\date{}

\newcommand{\R}{\mathbb{R}}
\newcommand{\Prob}{\mathbb{P}}
\newcommand{\one}{\boldsymbol{1}}
\newcommand{\Idea}{\paragraph{Idea.}}
\newcommand{\Esempio}{\paragraph{Esempio.}}
\newcommand{\Ricorda}{\paragraph{Da ricordare.}}
\newcommand{\Salta}{\paragraph{Cosa puoi saltare per ora.}}

\begin{document}
\maketitle

\paragraph{Come leggere questi appunti.}
Questa è una guida personale al secondo paper, da leggere dopo Schmidt. Non è
un testo da consegnare al relatore: è un percorso per arrivare alla riunione
sapendo raccontare il problema, capire i simboli principali e riconoscere la
parte più tecnica. Le figure sono schematiche, non dati reali. Quando scrivo
``$\lambda$ di Schmidt'' intendo il tail-dependence coefficient (TDC); quando
scrivo ``parametro $\lambda$ di Hua--Joe'' intendo invece il limite della
funzione lentamente variabile, che è una notazione diversa.

\section{Da Schmidt a Hua--Joe: quando $\lambda=0$ non basta}

In Schmidt la domanda di base è: se una variabile è già in una coda estrema,
quanto è probabile che anche l'altra sia estrema? Il limite di questa
probabilità è il TDC, per esempio $\lambda_L$ nella coda inferiore. Se il
risultato è positivo, abbiamo dipendenza asintotica; se è zero, parliamo di
indipendenza asintotica.

Il problema è che ``zero'' raccoglie moltissimi casi. Supponiamo che, per una
soglia rappresentata da $u\downarrow0$, due modelli abbiano probabilità
congiunte circa $u^{3/2}$ e $u^2$. Entrambe tendono a zero, quindi entrambe
hanno TDC nullo. Ma per $u=0.01$ otteniamo
\[
  u^{3/2}=10^{-3},\qquad u^2=10^{-4}.
\]
La prima probabilità è dieci volte più grande: in un problema di rischio non è
una differenza trascurabile. A soglie finite la probabilità è dunque non nulla;
la parola ``asintotica'' riguarda soltanto ciò che succede quando la soglia
viene spinta infinitamente lontano.

\Idea Hua e Joe propongono di conservare l'informazione sulla velocità con cui
la probabilità tende a zero. Chiamano questa velocità \emph{tail order} e
chiamano \emph{intermediate tail dependence} la regione tra dipendenza
asintotica e indipendenza completa.

\begin{figure}[ht]
\centering
\begin{tikzpicture}[x=0.9cm,y=0.85cm,font=\small]
  \draw[->] (0,0)--(8.2,0) node[right] {soglia verso l'estremo};
  \draw[->] (0,0)--(0,4.1) node[above] {probabilità congiunta};
  \draw[blue!70,thick] (0.3,3.4) .. controls (1.5,2.9) and (3.2,2.0) .. (7.8,0.45);
  \draw[red!70,thick] (0.3,3.4) .. controls (1.5,1.6) and (3.2,0.65) .. (7.8,0.05);
  \node[blue!70,align=left] at (5.2,2.0) {decadimento più lento\\ordine più piccolo};
  \node[red!70,align=left] at (5.1,0.65) {decadimento più rapido\\ordine più grande};
  \node[align=center] at (3.9,-0.75) {entrambe tendono a $0$:\\il solo TDC non le distingue};
\end{tikzpicture}
\caption{Due probabilità congiunte possono avere lo stesso limite nullo ma velocità diverse. Schema concettuale, non dati reali.}
\label{fig:decadimenti}
\end{figure}

\Ricorda Un ordine più piccolo significa, in generale, una discesa più lenta e
quindi più persistenza degli estremi. Il paper controlla anche il fattore
lentamente variabile che accompagna la potenza: l'esponente è il primo livello
di informazione, non sempre l'ultimo.

\section{Prima la storia dei dati: copula, marginali e diagonale}

Immagina una coppia di rendimenti giornalieri, ad esempio
$(r_{\mathrm{BMW}},r_{\mathrm{DB}})$. Una singola osservazione bivariata è
proprio questa coppia di numeri. La marginale BMW è invece il comportamento di
$r_{\mathrm{BMW}}$ preso da solo: la sua scala, la sua media e la sua coda. La
seconda marginale descrive Deutsche Bank. La copula non decide le unità di
misura: racconta come i due percentili si combinano.

Se BMW è al percentile $0.98$ e Deutsche Bank al percentile $0.97$, la coppia
si trova vicino all'angolo superiore destro del quadrato dei percentili,
indipendentemente dal fatto che i rendimenti abbiano volatilità diverse. È
questo posizionamento relativo che serve per le code.

Per marginali continue, la rappresentazione di Sklar è
\begin{equation}
 F(x_1,\ldots,x_d)=C\bigl(F_1(x_1),\ldots,F_d(x_d)\bigr),
 \label{eq:sklar}
\end{equation}
dove $d$ è il numero di componenti e $C$ è la copula. Se poniamo
$U_i=F_i(X_i)$, ogni $U_i$ è uniforme in $[0,1]$: abbiamo trasformato dati con
scale diverse in percentili confrontabili.

\begin{figure}[ht]
\centering
\begin{tikzpicture}[font=\scriptsize,>=Latex]
  \begin{scope}[xshift=0cm]
    \draw[->] (0,0)--(3.2,0) node[right] {BMW};
    \draw[->] (0,0)--(0,3.2) node[above] {DB};
    \foreach \x/\y in {0.5/0.7,0.8/0.5,1.0/1.2,1.3/1.4,1.5/1.1,1.8/2.0,2.1/1.9,2.4/2.5,2.8/2.3}{\fill[blue!65] (\x,\y) circle (1.5pt);}
    \node[align=center] at (1.5,-0.6) {spazio originale\\marginali diverse};
  \end{scope}
  \draw[->,thick] (3.5,1.5)--(4.5,1.5) node[midway,above] {percentili};
  \begin{scope}[xshift=5cm]
    \draw[->] (0,0)--(3.2,0) node[right] {$u=F_1(x_1)$};
    \draw[->] (0,0)--(0,3.2) node[above] {$v=F_2(x_2)$};
    \foreach \x/\y in {0.5/0.7,0.8/0.5,1.0/1.2,1.3/1.4,1.5/1.1,1.8/2.0,2.1/1.9,2.4/2.5,2.8/2.3}{\fill[blue!65] ({\x/3},{\y/3}) circle (1.5pt);}
    \node[align=center] at (1.5,-0.6) {percentili\\in $[0,1]^2$};
  \end{scope}
  \draw[->,thick] (8.5,1.5)--(9.5,1.5) node[midway,above] {copula};
  \begin{scope}[xshift=10cm]
    \draw[->] (0,0)--(3.2,0) node[right] {$u$};
    \draw[->] (0,0)--(0,3.2) node[above] {$v$};
    \draw[gray] (0,0) rectangle (3,3);
    \draw[dashed] (0,0)--(3,3);
    \foreach \x/\y in {0.5/0.7,0.8/0.5,1.0/1.2,1.3/1.4,1.5/1.1,1.8/2.0,2.1/1.9,2.4/2.5,2.8/2.3}{\fill[red!65] ({\x/3},{\y/3}) circle (1.5pt);}
    \node[align=center] at (1.5,-0.6) {legame\\tra percentili};
  \end{scope}
\end{tikzpicture}
\caption{Dal dato bivariato alla copula: prima si separano le marginali, poi si osserva come i percentili si combinano. Schema didattico.}
\label{fig:copula-pipeline}
\end{figure}

Per la lower tail guardiamo la diagonale vicino all'origine,
$C(u\one)=C(u,\ldots,u)$, con $u\downarrow0$. Per la upper tail si usa la
\emph{survival copula}:
\[
 \overline C(u_1,\ldots,u_d)=\Prob(U_1>u_1,\ldots,U_d>u_d),
\]
e la diagonale $\overline C((1-u)\one)$ con $u\downarrow0$. In dimensione due
$\overline C(u,v)=1-u-v+C(u,v)$; in dimensione maggiore interviene la formula
di inclusione--esclusione sulle copule marginali.

\begin{figure}[ht]
\centering
\begin{tikzpicture}[x=3.2cm,y=3.2cm,font=\small]
  \draw[->] (0,0)--(1.08,0) node[right] {$u$};
  \draw[->] (0,0)--(0,1.08) node[above] {$v$};
  \draw[gray] (0,0) rectangle (1,1);
  \draw[dashed,thick] (0,0)--(1,1) node[above right] {diagonale $u=v$};
  \fill[orange!25] (0,0) rectangle (0.27,0.27);
  \fill[blue!20] (0.73,0.73) rectangle (1,1);
  \node[align=center] at (0.14,0.14) {lower\\$C(u,u)$};
  \node[align=center] at (0.86,0.86) {upper\\$\overline C(1-u,1-u)$};
  \node[align=center] at (0.5,-0.22) {quadrato dei percentili};
\end{tikzpicture}
\caption{Le due regioni che si confrontano lungo la diagonale della copula. Schema concettuale, non dati reali.}
\label{fig:copula-diagonal}
\end{figure}

\Salta per ora le formule complete di inclusione--esclusione: è sufficiente
ricordare che la survival copula misura probabilità congiunte di superamento,
mentre $C$ misura probabilità congiunte di non superamento.

\section{Tail order $\kappa$: prima l'idea, poi la formula}

Hua e Joe descrivono una diagonale come una potenza moltiplicata per una
correzione lenta:
\begin{equation}
 C(u\one)\sim u^{\kappa_L(C)}\ell_L(u),
 \qquad u\downarrow0,
 \label{eq:kappa-lower}
\end{equation}
e, per la coda superiore,
\begin{equation}
 \overline C((1-u)\one)\sim u^{\kappa_U(C)}\ell_U(u),
 \qquad u\downarrow0.
 \label{eq:kappa-upper}
\end{equation}

Leggiamo la prima formula lentamente. $u$ è la distanza dalla coda; $u^\kappa$
è la velocità principale; $\ell(u)$ è una funzione lentamente variabile, per
esempio una potenza di $-\log u$. Il simbolo $\sim$ significa che il rapporto
tra i due membri tende a $1$. Il paper giustifica questa scrittura con la
teoria della variazione regolare:
\[
 \lim_{x\to\infty}\frac{g(tx)}{g(x)}=t^\alpha,qquad t>0.
\]

\begin{center}
\begin{tabular}{p{2.7cm}p{3.5cm}p{5.1cm}}
\toprule
Ordine & Situazione & Come dirlo a voce\\
\midrule
$\kappa=1$ & dipendenza di coda, se il fattore lento tende a una costante positiva & La probabilità congiunta è dell'ordine di $u$; qui compare il TDC positivo.\\
$1<\kappa<d$ & intermediate tail dependence, sotto dipendenza positiva & La probabilità va a zero, ma più lentamente dell'indipendenza completa.\\
$\kappa=d$ & riferimento di indipendenza & La diagonale ha l'ordine di $u^d$; con la copula indipendente è proprio $u^d$.\\
\bottomrule
\end{tabular}
\end{center}

\paragraph{Due lambda diverse.}
In Schmidt $\lambda_L$ e $\lambda_U$ sono limiti di probabilità condizionate:
sono i TDC. In Hua--Joe, invece, il simbolo $\lambda_L(C)$ o
$\lambda_U(C)$ viene usato per il limite di $\ell(u)$, quando esiste. Se
$\kappa=1$ e $\ell(u)\to h>0$, questo $h$ coincide con il TDC nella
convenzione della coda considerata. Se $\kappa>1$, il TDC è zero anche se
guardare $\kappa$ è informativo. Non bisogna quindi scambiare i due ruoli.

\Ricorda Più piccolo è $\kappa$, più lento è il decadimento di $u^\kappa$ e
maggiore è la persistenza della dipendenza estrema. Una copula non può avere
$\kappa<1$ nel caso ordinario perché $C(u\one)\leq u$. Per una copula
indipendente in dimensione $d$, invece, $\kappa=d$.

\Esempio In dimensione due e per $u=0.01$:
\[
 u^1=10^{-2},\qquad u^{3/2}=10^{-3},\qquad u^2=10^{-4}.
\]
Il valore $3/2$ è davvero intermedio: non produce il limite positivo della
dipendenza asintotica, ma non è neppure il decadimento rapido $u^2$.

Il paper collega questa notazione a Ledford--Tawn, che scrivono una probabilità
congiunta come $\ell(r)r^{-1/\eta}$, con $1/2\leq\eta\leq1$. La relazione è
$\kappa=1/\eta$. È un collegamento utile, ma non serve farne il centro della
prima lettura.

\section{Percorsi non diagonali: la tail order function}

La diagonale confronta soglie uguali in tutte le componenti. Ma in pratica
una variabile può essere estremamente lontana dalla propria media mentre
un'altra è solo moderatamente estrema. Per questo gli autori assegnano pesi
$w=(w_1,\ldots,w_d)$ e studiano
\begin{equation}
 b(w;C,\kappa)=\lim_{u\downarrow0}
 \frac{C(uw_1,\ldots,uw_d)}{u^\kappa\ell(u)},
 \label{eq:order-function}
\end{equation}
quando il limite esiste. La funzione upper $b^*$ sostituisce il numeratore con
$\overline C(1-uw_1,\ldots,1-uw_d)$.

\Idea $b(w)$ descrive la forma della dipendenza vicino all'angolo, non solo il
numero letto sulla diagonale. Se prendiamo $w=\one$ torniamo alla definizione
di $\kappa$. Gli autori mostrano anche l'omogeneità $b(tw)=t^\kappa b(w)$:
spostare tutti i pesi nella stessa proporzione cambia la scala con l'ordine
stesso.

\Salta per ora le dimostrazioni e le derivate della funzione $b$. Il concetto
da conservare è che due copule con lo stesso $\kappa$ possono ancora avere
forme diverse lungo percorsi non diagonali.

\section{La copula gaussiana: esempio di intermediate tail dependence}

La copula gaussiana d-dimensionale è costruita trasformando i percentili in
quantili normali e usando una matrice di correlazione $\Sigma$:
\begin{equation}
 C_d^\Phi(u_1,\ldots,u_d)=
 \Phi_d\!\left(\Phi^{-1}(u_1),\ldots,\Phi^{-1}(u_d);\Sigma\right).
 \label{eq:gaussian-copula}
\end{equation}

Nel caso bivariato con correlazione positiva $\rho$ gli autori ottengono
\begin{equation}
 \kappa=\frac{2}{1+\rho},
 \qquad
 \ell(u)=(-\log u)^{-\rho/(1+\rho)}.
 \label{eq:gaussian-order}
\end{equation}
Se $\rho=0.5$, allora $\kappa=4/3$. La copula è ancora asintoticamente
indipendente nel senso di Schmidt ($\lambda=0$), ma la correlazione positiva ha
rallentato il decadimento rispetto a $\kappa=2$ dell'indipendenza. In
dimensione $d$ con correlazione equicorrelata:
\[
 \kappa=\frac{d}{1+(d-1)\rho}.
\]

\paragraph{Come raccontarlo al relatore.}
``Una gaussiana con correlazione positiva non crea un TDC positivo, però la
probabilità di estremi congiunti diminuisce più lentamente. Quindi il TDC
classico la mette nella casella indipendenza asintotica, mentre il tail order
la colloca nella zona intermedia.'' Questa è esattamente l'informazione che
Schmidt, da solo, non separava.

\section{Archimedean copulas: il fattore latente prima dei simboli}

Una copula archimedea può essere vista come un modello con un fattore comune
non osservato. Condizionatamente a quel fattore, le componenti sono
indipendenti; quando il fattore assume valori particolari, però, tutte le
componenti vengono spinte verso lo stesso bordo. È un modo intuitivo per
generare dipendenza.

Matematicamente il paper usa la trasformata di Laplace di una variabile
$Y\geq0$:
\begin{equation}
 C_\psi(u_1,\ldots,u_d)=
 \psi\!\left(\psi^{-1}(u_1)+\cdots+\psi^{-1}(u_d)\right),
 \qquad \psi(s)=\mathbb E[e^{-sY}].
 \label{eq:arch-lt}
\end{equation}

Se la variabile latente ha una coda superiore pesante, c'è più probabilità di
spingere simultaneamente le componenti verso il bordo superiore. Se ha molta
densità vicino a zero, si rafforza la lower tail. Questa è la ragione per cui
le copule archimedee possono essere asimmetriche tra le due code.

\subsection{Il momento massimo e il tail order}
Gli autori riassumono la pesantezza della coda di $Y$ con
\begin{equation}
 M_\psi=\sup\{m\geq0:\mathbb E(Y^m)<\infty\}.
 \label{eq:max-moment}
\end{equation}
Un valore piccolo di $M_\psi$ indica che già potenze basse di $Y$ hanno
momento infinito, quindi la coda è pesante. Se $k<M_\psi<k+1$, con
$k\in\{1,\ldots,d-1\}$, la Proposizione 4 dà un ordine upper
$\kappa_U=M_\psi$: è intermediate tail dependence. Se tutte le derivate
necessarie di $\psi$ sono finite, l'ordine può arrivare a $d$; se
$\psi'(0)=-\infty$ e $0<M_\psi<1$, l'ordine è $1$ e nel caso bivariato
\begin{equation}
 \lambda_U^{\text{(TDC)}}=2-2^{M_\psi}.
 \label{eq:arch-upper-lambda}
\end{equation}
Qui ho scritto esplicitamente ``TDC'' per non confonderlo con il parametro
$\lambda_U(C)=\lim\ell_U(u)$ della Definizione 2.

\Salta per ora la dimostrazione con tutte le derivate della LT. Per la prima
comprensione basta sapere che una caratteristica univariata della variabile
latente determina l'ordine di una dipendenza multivariata.

\subsection{Lower tail e velocità della LT}
Per la coda inferiore il paper studia il comportamento di $\psi(s)$ quando
$s\to\infty$, ad esempio
\[
 \psi(s)\sim a_1s^q\exp\{-a_2s^r\}.
\]
Quando $0<r<1$, l'ordine lower ottenuto è $\kappa_L=dr$; per $r=0$ l'ordine
è $1$, per $r=1$ è $d$. La lettura intuitiva è che una discesa più lenta della
LT per $s$ grande segnala più massa vicino a zero nella variabile latente e
quindi più dipendenza nella lower tail.

\section{La famiglia inverse-Gamma: un esempio di asimmetria progettata}

Hua e Joe non si limitano a classificare famiglie già note: propongono una
famiglia a un parametro per ottenere pattern di coda che Gaussiane, $t$ e
molte archimedean standard non rappresentano bene. Se
$X\sim\operatorname{Gamma}(\alpha,1)$ e $Y=X^{-1}$, la sua LT è
\begin{equation}
 \psi(s;\alpha)=\frac{2}{\Gamma(\alpha)}s^{\alpha/2}
 K_\alpha(2\sqrt{s}),\qquad \alpha>0,
 \label{eq:inverse-gamma-lt}
\end{equation}
dove $K_\alpha$ è la funzione di Bessel modificata del secondo tipo.

Per $\alpha$ non intero, la Proposizione 7 dà
\begin{equation}
 \kappa_U=\max\{1,\min\{\alpha,d\}\},
 \qquad \kappa_L=\sqrt d.
 \label{eq:inverse-gamma-orders}
\end{equation}
In dimensione due, $\kappa_L=\sqrt2$ e
$\kappa_U=\max\{1,\min\{\alpha,2\}\}$. Se $0<\alpha<\sqrt2$,
$\kappa_U<\kappa_L$: l'upper order è più piccolo, quindi la upper tail è più
persistente. Se $\alpha>\sqrt2$ accade il contrario. Numericamente, la
concordanza centrale cresce verso la copula comonotona quando $\alpha$ scende
a zero e tende all'indipendenza quando $\alpha$ cresce.

\begin{center}
\begin{tabular}{p{3.2cm}p{2.7cm}p{2.7cm}p{4.2cm}}
\toprule
Regione di $\alpha$ & $\kappa_U$ & $\kappa_L$ & Lettura in parole\\
\midrule
$\alpha\downarrow0$ & vicino a $1$ & $\sqrt d$ & upper tail molto persistente.\\
$\alpha=\sqrt2$ in $d=2$ & $\sqrt2$ & $\sqrt2$ & soglia di simmetria degli ordini.\\
$\alpha\to\infty$ & $d$ & $\sqrt d$ per ogni valore finito & la concordanza tende all'indipendenza; il limite centrale e il limite asintotico non vanno confusi.\\
\bottomrule
\end{tabular}
\end{center}

\paragraph{Perché è utile.}
La famiglia permette di cambiare il lato della distribuzione più persistente
senza dover cambiare completamente la costruzione. È un esempio concreto del
messaggio del paper: tail order e simmetria di riflessione possono essere
progettati separatamente dalla dipendenza centrale.

\section{Miscele max-infinitely divisible: combinare meccanismi}

La seconda costruzione importante parte da una copula max-infinitely divisible
$K$ e da una LT $\psi$:
\begin{equation}
 C(u_1,\ldots,u_d)=
 \psi\!\left[-\log K\!\left(e^{-\psi^{-1}(u_1)},\ldots,
 e^{-\psi^{-1}(u_d)}\right)\right].
 \label{eq:maxid-mixture}
\end{equation}

La formula dice, in sostanza, che si inserisce una struttura $K$ dentro il
meccanismo archimedeo della LT. Se $K$ porta già un certo ordine e $\psi$ ha
una certa pesantezza, il risultato combina i due effetti. Le Proposizioni 8 e
9 mostrano condizioni in cui l'ordine upper resta quello di $C_\psi$ oppure
diventa un nuovo ordine intermedio determinato da entrambi.

\Idea È come avere due manopole: una per il meccanismo di dipendenza della
copula $K$, una per il comportamento della variabile latente. In dimensione
due questo produce molte candidate utili come linking copulas; in dimensione
maggiore le sole archimedean copulas sono exchangeable, quindi le vine copulas
offrono un modo per assemblare blocchi bivariati più flessibili.

\Salta per ora i coefficienti di inclusione--esclusione nelle Proposizioni 8--9.
Il punto da ricordare è che le miscele permettono upper/lower tail diverse e
una struttura più ricca di quella di una singola famiglia simmetrica.

\section{Reflection symmetry: capire quale coda domina}

Riflettere una copula significa sostituire $(U_1,\ldots,U_d)$ con
$(1-U_1,\ldots,1-U_d)$. Se la copula resta uguale, c'è reflection symmetry.
Se non resta uguale, confrontiamo gli ordini:
\begin{itemize}[nosep]
  \item $\kappa_L>\kappa_U$: l'upper tail ha ordine più piccolo e quindi
  decresce più lentamente; l'asimmetria è verso l'upper tail;
  \item $\kappa_L<\kappa_U$: la lower tail è più persistente;
  \item $\kappa_L=\kappa_U$ non garantisce simmetria: i fattori $\ell$ o le
  funzioni di ordine possono ancora essere diversi.
\end{itemize}

\begin{figure}[ht]
\centering
\begin{tikzpicture}[x=1.3cm,y=0.75cm,font=\small]
  \draw[->] (0,0)--(5.9,0) node[right] {modello};
  \draw[->] (0,0)--(0,3.8) node[above] {persistenza della coda};
  \draw[gray!60,dashed] (0.2,1.0)--(5.6,1.0);
  \node[gray,anchor=west] at (0.25,0.92) {indipendenza: $\kappa=d$};
  \draw[blue!70,very thick] (0.5,3.0)--(2.0,2.4)--(3.5,1.7)--(5.3,1.0);
  \draw[red!70,very thick] (0.5,2.0)--(2.0,2.0)--(3.5,1.5)--(5.3,1.0);
  \node[blue!70,align=left] at (1.1,3.35) {ordine più piccolo:\\più persistenza};
  \node[red!70,align=left] at (3.5,2.2) {ordine più grande:\\decadimento più rapido};
  \node[align=center] at (2.9,-0.65) {$\kappa=1$ è la coda dipendente;\\$1<\kappa<d$ è la zona intermedia};
\end{tikzpicture}
\caption{Schema qualitativo della scala di tail order. Le ordinate non sono valori stimati.}
\label{fig:classification}
\end{figure}

\section{Conclusione: cosa cambia nella mia mappa mentale}

\subsection*{Schmidt e Hua--Joe in una frase ciascuno}
Schmidt fornisce il primo linguaggio operativo: TDC, copule, stimatori e VaR.
Hua--Joe prende la casella ``TDC uguale a zero'' e la apre: due modelli
asintoticamente indipendenti possono avere velocità di decadimento molto
diverse, e questa velocità può essere asimmetrica tra upper e lower tail.

\subsection*{Cosa devo saper dire al relatore}
Potrei dire: ``Il TDC dice se la probabilità condizionata estrema resta
positiva. Se è zero, Hua--Joe guarda la probabilità congiunta e la scrive come
$u^\kappa\ell(u)$. L'esponente $\kappa$ misura la velocità: $1$ è la coda
dipendente, $d$ è l'indipendenza e i valori intermedi descrivono una
dipendenza estrema che tende comunque a zero. La copula gaussiana con $\rho>0$
è l'esempio semplice: TDC nullo ma ordine intermedio.''

\subsection*{Cosa non devo ancora saper fare}
Non devo ancora ricostruire le dimostrazioni con trasformate di Laplace,
derivate di ordine elevato o funzioni di Bessel. Devo però sapere che quelle
tecniche collegano una proprietà univariata della variabile latente alla coda
della copula. Non devo neppure avere già scelto una famiglia per la tesi:
questo paper serve prima a capire quali comportamenti sarebbe possibile
rappresentare.

\subsection*{Rilettura consigliata}
Per un secondo passaggio leggerei le Definizioni 1--3 alle pp.~1455--1457,
poi l'esempio gaussiano a p.~1456. Solo dopo tornerei alle Proposizioni 4--7
(pp.~1460--1462), guardando soprattutto la direzione delle disuguaglianze e
non ogni passaggio della prova. Infine leggerei le Proposizioni 8--9 a p.~1462
per capire la motivazione delle miscele max-id.

\subsection*{Parole chiave}
tail order, intermediate tail dependence, TDC, $\lambda_L$, $\lambda_U$,
$\kappa_L$, $\kappa_U$, variazione regolare, funzione lentamente variabile,
diagonal section, survival copula, Archimedean copula, Laplace transform,
maximal moment, inverse-Gamma, max-infinitely divisible, reflection symmetry,
tail asymmetry, vine copula.

\subsection*{Legenda finale della notazione e pronuncia}
\begin{description}[leftmargin=3.5cm,style=nextline]
  \item[$\lambda_L,\lambda_U$ in Schmidt] TDC lower e upper, cioè limiti di
  probabilità condizionate. Lambda si legge ``lambda''.
  \item[$\lambda_L(C),\lambda_U(C)$ in Hua--Joe] limiti del fattore lento
  $\ell(u)$ nella Definizione 2; coincidono con il TDC solo nel caso
  $\kappa=1$ e limite positivo.
  \item[$\kappa_L,\kappa_U$] tail order, esponente della velocità; kappa è la
  lettera greca $\kappa$ e si legge ``cappa''.
  \item[$k$] indice intero usato nelle prove per individuare l'intervallo
  $k<M_\psi<k+1$; è la lettera latina ``k'', non il parametro $\kappa$.
  \item[$\theta,\alpha$] parametri di famiglie; in Hua--Joe $\alpha$ è il
  parametro inverse-Gamma e $M_\psi=\alpha$ per quella famiglia.
  \item[$\rho$] correlazione della copula gaussiana bivariata.
  \item[$C,\overline C$] copula ordinaria e survival copula.
  \item[$u,v$] livelli percentili in $[0,1]$; $u\downarrow0$ indica la lower
  tail, mentre $1-u\uparrow1$ indica la upper tail.
  \item[$\one$] vettore $(1,\ldots,1)$; $C(u\one)$ è la diagonale.
  \item[$\ell(u)$] funzione lentamente variabile che corregge $u^\kappa$.
  \item[$\lim$, $\sim$] limite e equivalenza asintotica, cioè rapporto tendente
  a $1$.
  \item[LT] trasformata di Laplace, $\psi(s)=\mathbb E[e^{-sY}]$.
  \item[$M_\psi$] massimo ordine di momento finito della variabile positiva
  $Y$ associata alla LT.
\end{description}

\subsection*{Domande per il relatore}
\begin{description}[leftmargin=3.1cm,style=nextline]
  \item[Concettuali] Nel nostro lavoro conviene usare $\kappa$ anche quando
  il TDC di Schmidt è zero? Come distinguere intermediate tail dependence da
  near independence? Quanto è importante il fattore lentamente variabile?
  \item[Metodologiche] Esistono stimatori affidabili di $\kappa$ e $b(w)$ per
  campioni finanziari realistici? Come scegliere tra copula gaussiana, $t$,
  archimedea e vine quando le code sono asimmetriche?
  \item[Empiriche] Come verificare la reflection asymmetry nei rendimenti? Ha
  senso stimare separatamente lower e upper tail nel VaR invece di usare una
  copula $t$ simmetrica?
  \item[Collegamento ai paper] Schmidt offre il quadro operativo TDC--stima--
  VaR; Hua--Joe approfondisce la classificazione quando il TDC è nullo.
  Durante et al. sembra concentrarsi sulla diagonal section e sugli ordini di
  coda in una direzione più geometrica. Chiedere se il relatore vede questi
  lavori come una sequenza naturale o come approcci alternativi.
\end{description}

\end{document}
